\begin{abstract}
Although large-scale unsupervised language models (LMs) acquire extensive knowledge about the world and display some reasoning capabilities, exerting fine-grained control over their outputs remains challenging because their training lacks supervision. To address this limitation, current approaches leverage human-annotated comparisons of model outputs and subsequently fine-tune the LM to better reflect these preferences, frequently utilizing reinforcement learning from human feedback (RLHF). Nonetheless, RLHF typically involves a two-stage and sometimes unstable process: first training a reward model to encapsulate human preferences, followed by optimizing the original LM via reinforcement learning to increase this proxy reward while minimizing divergence from the base model. In this work, we propose a novel reparameterization of the RLHF reward model, which allows the corresponding optimal policy to be directly computed in closed form. This reformulation lets us address the RLHF objective using a straightforward classification loss alone. The new method, termed \textit{Direct Preference Optimization} (DPO), offers stability, strong empirical performance, and computational efficiency, removing the requirements for LM sampling during fine-tuning or extensive hyperparameter adjustment. Experimental results demonstrate that DPO matches or surpasses prior techniques in aligning LMs to human preferences. Notably, DPO-based fine-tuning outperforms RLHF using PPO for generation sentiment control, and achieves comparable or superior quality in summarization and single-turn dialogue—all while offering a much simpler implementation and training process.
\end{abstract}

\section{Introduction}
Large unsupervised language models (LMs), when trained on massive corpora, demonstrate an array of unexpected capabilities~\citep{chowdhery2022palm, brown2020language, touvron2023llama, bubeck2023sparks}. Yet, the underlying training data reflects human contributions with heterogeneous intentions, priorities, and competencies. Certain aims or expertise captured in the data may not be ideal for models to replicate. For instance, while it is advantageous for an AI programming assistant to \textit{recognize} frequent coding errors so it can address them, we nonetheless prefer that its code generation be influenced predominantly by the relatively rare but superior programming exemplars embedded in its training dataset. Similarly, though it may be beneficial for a language model to \textit{identify} a prevalent misconception held by half of the population, we would not want the model to reproduce that misconception in half its responses to relevant queries. Thus, effectively \emph{choosing} how the model responds and acts—drawing selectively from its extensive \textit{knowledge and capabilities}—is fundamental for constructing AI systems that are reliable, high-performing, and manageable~\citep{ouyang2022training}. While prevalent methodologies align LMs with human preferences via reinforcement learning (RL), we demonstrate that the RL-centric objectives of current methods can be satisfied precisely using a straightforward binary cross-entropy objective, thereby significantly streamlining the process of preference alignment.

Broadly, current techniques instill target behaviors in language models by leveraging carefully constructed datasets of human preferences, which encapsulate notions of safe and useful conduct. This stage of preference alignment succeeds the large-scale, unsupervised pre-training phase on extensive textual corpora. Although a naive approach would employ supervised fine-tuning on human-annotated, high-quality outputs, the predominant and most effective strategies are based on reinforcement learning from human or AI feedback (RLHF/RLAIF; \citep{christiano2017deep,bai2022constitutional}). In RLHF, a reward model is trained on datasets reflecting human judgments, and subsequently, RL is used to adapt the LM’s policy to favor outputs deemed highly rewarding while constraining deviation from the initial model. Despite the considerable advancements made possible by RLHF—such as impressive capabilities in conversation and code synthesis—the methodology remains complex compared to straightforward supervised learning, demanding the training of several LMs and frequent in-the-loop policy sampling, leading to heightened computational resource requirements.

This work demonstrates a method for optimizing language models according to human preferences without relying on explicit reward modeling or reinforcement learning frameworks. We introduce \textit{{Direct Preference Optimization} (DPO)}, an approach that inherently seeks to optimize the same objective pursued by conventional RLHF techniques—maximizing reward under a KL-divergence regularization—while offering simplicity in implementation and ease of training. DPO achieves this by adjusting the relative log probabilities between preferred and non-preferred outputs, moderated by an adaptive, example-specific importance weight. This weighting mechanism mitigates the model degeneration issues associated with a straightforward application of a probability ratio objective. Similar to previous approaches, DPO employs a formal preference model (for instance, the Bradley-Terry model; \cite{bradley1952rankanalysis}) to assess how closely a reward function aligns with observed preference data. Nonetheless, rather than leveraging this preference model first to train a reward model and subsequently to train a policy to maximize rewards, DPO directly expresses the preference loss as a function of the policy through a variable substitution. Consequently, when provided with a dataset of human-annotated preferences for model responses, DPO allows direct policy optimization via a binary cross-entropy loss, thereby yielding the optimal policy corresponding to a reward function implicitly fitted to the preferences.

The central contribution of our study is Direct Preference Optimization (DPO), an uncomplicated, reinforcement learning-free approach for training language models based on preference data. Our empirical evaluations indicate that DPO matches or surpasses the performance of prevailing techniques, such as PPO-based RLHF, for preference learning tasks—spanning sentiment control, summarization, and conversational applications—utilizing language models of up to 6B parameters.

\section{Related Work}

As self-supervised language models have grown in size, they have demonstrated the capability to perform certain tasks in a zero-shot setting \citep{radford2019language}, or using few-shot prompting strategies \citep{gpt3,megatron,chowdhery2022palm}. Nevertheless, their effectiveness on downstream applications and consistency with user objectives can be considerably enhanced by fine-tuning on datasets comprised of explicit instructions and corresponding human-generated completions \citep{mishra-etal-2022-cross,sanh2022multitask,chung2022scaling,thoppilan2022lamda}. The process known as `instruction-tuning’ equips LLMs with the capacity to handle instructions that extend beyond those seen during fine-tuning, thereby boosting their general applicability \citep{chung2022scaling}. Although instruction-tuning has proven successful, it is often more feasible to gather \textit{relative} human assessments of response quality compared to assembling expert-curated demonstrations. For this reason, subsequent research has sought to fine-tune LLMs using datasets reflecting human preferences, which has yielded improvements in areas such as translation \citep{kreutzer-etal-2018-reliability}, summarization \citep{stiennon2022learning,ziegler2020finetuning}, narrative generation \citep{ziegler2020finetuning}, and instruction following \citep{ouyang2022training,ramamurthy2023is}. These approaches typically start by training a neural network to predict preference scores aligned with collected preference data, often modeled with the Bradley-Terry framework \citep{bradley1952rankanalysis}. Subsequently, the language model is adjusted to maximize these predicted rewards via reinforcement learning techniques, such as REINFORCE \citep{williams1992reinforce}, proximal policy optimization (PPO; \cite{schulman2017proximal}), or related methodologies \citep{ramamurthy2023is}. Another closely linked research direction utilizes LLMs, already refined for instruction-following via human feedback, to autonomously produce further synthetic preference annotations focusing on particular qualities like safety or harmlessness \citep{bai2022constitutional}, relying primarily on minimal human input in the form of evaluative text rubrics for annotation guidance. Collectively, these strategies represent the intersection of research on reinforcement learning-based language model training for varied goals~\citep{Ranzato2015SequenceLT,paulus2018a,wu2018learning} and on foundational techniques for preference-based learning from human feedback \citep{christiano2017deep,kupcsik2018learning}. Despite the promise of using relative preference data, refining large language models via reinforcement learning is still practically challenging; this paper presents a theoretically grounded alternative for optimizing over relative preferences without the use of RL.

Beyond linguistic applications, learning policies from preferences has been explored within both bandit and reinforcement learning frameworks, with a range of methodologies proposed. In particular, contextual bandit approaches that utilize preferences or rankings over actions instead of direct rewards are known as contextual dueling bandits (CDB; \cite{yue2012karmed,dudik2015contextual}). When absolute rewards are unavailable, theoretical analyses of CDBs often replace the traditional optimal policy concept with the idea of a \textit{von Neumann winner}, defined as a policy whose probability of outperforming any alternative policy meets or exceeds 50\% \citep{dudik2015contextual}. Notably, preference feedback in CDBs is typically acquired in an online manner, whereas, in human preference learning scenarios, algorithms generally operate on a finite, static collection of preference-annotated action pairs \citep{yan2022human}. Preference-based RL (PbRL) similarly aims to learn from binary preferences, which are sampled from some unknown 'scoring' function as opposed to observed rewards \citep{BusaFekete2014,ruiz2023dueling}. Numerous PbRL algorithms have been devised, with several leveraging off-policy preference data; nonetheless, these methods frequently require the explicit estimation of the hidden scoring (reward) function, followed by subsequent policy optimization \citep{jain2013learning,BusaFekete2014,christiano2017deep,sadigh2017active,kupcsik2018learning}. In contrast, our contribution is a unified policy learning framework that directly trains the policy to fulfill specified preferences, without explicitly modeling the underlying reward.

\section{Preliminaries}\label{section:prelims}

Here, we revisit the RLHF process outlined by \citeauthor{ziegler2020finetuning} (with subsequent extensions in \citep{stiennon2022learning, bai2022training, ouyang2022training}). This framework generally comprises three stages: (1) supervised fine-tuning (SFT); (2) collecting preference data and reward model training; and (3) reinforcement learning-based policy optimization.

\textbf{SFT}: The RLHF workflow starts by adapting a pre-trained language model to downstream tasks (such as dialogue or summarization) via supervised learning on curated datasets, producing an initial model $\pi^\text{SFT}$.

\textbf{Reward Modelling Phase}: Next, prompts $x$ are fed into the SFT-tuned model to generate answer pairs $(y_1, y_2)\sim \pi^\text{SFT}(y \mid x)$. Human annotators are then tasked with indicating their preference between these outputs, denoted by $y_w\succ y_l \mid x$, where $y_w$ is the selected response and $y_l$ is the less favored one from the pair $(y_1, y_2)$. While the true reward signal $r^*(y, x)$ responsible for these preferences remains hidden, the literature offers several models to describe preference generation; the Bradley-Terry (BT) \cite{bradley1952rankanalysis} model is commonly employed (though the broader Plackett-Luce models \citep{plackett1975analysis, luce2012individual} can also be applied when rankings over more than two answers are available). Under the BT framework, the probability distribution $p^*$ expressing human preferences is given by:

\begin{equation}\label{eq:bradley-terry}
    p^*(y_1\succ y_2 \mid x)=\frac{\exp\left(r^*(x, y_1)\right)}{\exp\left(r^*(x, y_1)\right) + \exp\left(r^*(x, y_2)\right)}.
\end{equation}

Given a fixed dataset of pairwise comparisons $\mathcal{D} = \bigl\{x^{(i)}, y_w^{(i)}, y_l^{(i)}\bigr\}_{i=1}^N$ generated according to $p^*$, we introduce a parametrized reward function $r_{\phi}(x, y)$. The parameter estimation is performed through maximum likelihood. When this is viewed as a binary classification task, the negative log-likelihood objective can be written as follows:

\begin{equation}\label{eq:reward_model}
    \mathcal{L}_R(r_{\phi}, \mathcal{D}) = -\mathbb{E}_{(x, y_w, y_l)\sim \mathcal{D}}\bigl[\log \sigma(r_{\phi}(x, y_w)- r_{\phi}(x, y_l))\bigr]
\end{equation}

Here, $\sigma$ denotes the sigmoid function. In large language model (LM) contexts, the reward network $r_{\phi}(x, y)$ is commonly initialized with the SFT model $\pi^\text{SFT}(y \mid x)$, augmented by an additional linear head that outputs a scalar reward estimate based on the last hidden representation~\cite{ziegler2020finetuning}. To achieve more stable reward predictions, previous work employs a normalization constraint such that $\mathbb{E}_{x,y\sim \mathcal{D}}\left[r_\phi(x, y)\right] = 0$ holds for all $x$.

\textbf{RL Fine-Tuning Phase}: In the subsequent RL optimization step, the obtained reward model is used to guide the LM updates. In line with established practice~\citep{jaques2017sequence, jaques2020human}, the following regularized objective is considered:

\begin{equation}\label{eq:RL}
\max_{\pi_{\theta}}  \mathbb{E}_{x\sim \mathcal{D}, y\sim \pi_{\theta}(y \mid x)}\bigl[r_{\phi}(x, y)\bigr] - \beta\mathbb{D}_{\textrm{KL}}\bigl[\pi_{\theta}(y\mid x)\mid \mid \pi_\text{ref}(y\mid x)\bigr],
\end{equation}

with $\beta$ governing how strongly the updated policy is encouraged to remain close to the base policy $\pi_\text{ref}$, typically set to the initial SFT policy $\pi^\text{SFT}$. 
Practically, $\pi_\theta$ is initialized from $\pi^\text{SFT}$. This KL penalty serves a dual purpose: constraining the policy to stay within the reliable domain of the reward model and ensuring response variety, thus avoiding degenerate behavior such as mode collapse. As the language space is discrete, this objective is not directly differentiable and is hence tackled via reinforcement learning techniques. Commonly, the reward is defined as ${r(x, y) = r_{\phi}(x, y) -\beta (\log \pi_{\theta}(y\mid x) - \log \pi_\text{ref}(y\mid x))}$ and the policy is optimized using PPO~\cite{schulman2017proximal}, following the approach of \cite{ziegler2020finetuning, stiennon2022learning, bai2022training, ouyang2022training}.

\section{Direct Preference Optimization}\label{sec:DPO}

Recognizing the complexity inherent in applying standard RL methods to the scale of contemporary language model fine-tuning, we propose a streamlined preference-based policy optimization framework. In contrast to established RLHF pipelines that decouple preference modeling and RL optimization, our approach exploits a particular structure in the reward parameterization, allowing us to derive the associated optimal policy in closed form—eliminating the need for an explicit RL phase. We show in detail that by utilizing an analytic mapping between reward functions and corresponding optimal policies, we can convert objectives for fitting reward models into objectives for training policies directly. This variable substitution bypasses explicit reward fitting but remains consistent with the human preference modeling assumptions, such as the Bradley-Terry framework. Effectively, this makes the policy network simultaneously represent both the generative language model and an implicit reward estimator.

\textbf{Deriving the DPO objective.} Consider the standard RL objective as formalized in Eq.~\ref{eq:RL}, employing a generic reward function $r$. Building on the methodology of prior studies~\citep{peters2007reinforcement, peng2019advantage, korbak2022reinforcement, go2023aligning}, it can readily be established that the solution which optimizes the KL-regularized reward maximization problem in Eq.~\ref{eq:RL} has the form:

\begin{equation}\label{eq:op_policy}
    \pi_r(y\mid x) = \frac{1}{Z(x)}\pi_\text{ref}(y\mid x)\exp\left(\frac{1}{\beta}r(x, y)\right),
\end{equation}

where the partition function is defined as $Z(x) =\sum_{y}\pi_\text{ref}(y\mid x)\exp\left(\frac{1}{\beta}r(x, y)\right)$. A detailed proof can be found in Appendix \ref{app:derivation1}. Notably, even when substituting the reward function $r$ with its maximum likelihood estimate $r_\phi$ as an approximation of the true reward $r^*$, calculating $Z(x)$ remains computationally demanding~\citep{korbak2022reinforcement, go2023aligning}. This computational bottleneck makes direct use of this formulation impractical. Nevertheless, by manipulating Eq.~\ref{eq:op_policy}, one can rewrite the reward function in terms of the optimal policy $\pi_r$, the reference policy $\pi_\text{ref}$, and the as-yet undetermined partition function $Z(\cdot)$. Specifically, applying the logarithm to both sides of Eq.~\ref{eq:op_policy} and rearranging terms yields:

\begin{equation}\label{eq:main_eq}
    r(x,y) =\beta \log \frac{\pi_r(y\mid x)}{\pi_\text{ref}(y\mid x)} + \beta \log Z(x).
\end{equation}

This transformation applies equally to the ground-truth reward $r^*$ and its corresponding optimal policy $\pi^*$. Crucially, the Bradley-Terry model only requires knowledge of the difference in rewards for two candidate outputs, that is, ${p^*(y_1 \succ y_2 \mid x) = \sigma(r^*(x, y_1) - r^*(x, y_2))}$. Inserting the representation of $r^*(x, y)$ from Eq.~\ref{eq:main_eq} into the Bradley-Terry preference model given by Eq.~\ref{eq:bradley-terry}, the partition function terms cancel out. As a result, the probability of human preference can be described entirely using the optimal policy $\pi^*$ and the reference policy $\pi_\text{ref}$. This shows that under the Bradley-Terry model, the RLHF-optimal policy $\pi^*$ satisfies:

\begin{equation}\label{eq:objective}
    p^*(y_1\succ y_2 \mid x)=\frac{1}{1 + \exp\left(\beta \log \frac{\pi^*(y_2\mid x)}{\pi_\text{ref}(y_2\mid x)} - \beta \log \frac{\pi^*(y_1\mid x)}{\pi_\text{ref}(y_1\mid x)}\right)}
\end{equation}

A full derivation is provided in Appendix~\ref{app:derivation2}. While Eq.~\ref{eq:objective} utilizes the Bradley-Terry model, a parallel approach can be used to derive analogous expressions for more general Plackett-Luce models~\citep{plackett1975analysis, luce2012individual}, as detailed in Appendix~\ref{app:plackett_luce_models}.

Having recast the probability of human preferences in terms of the optimal policy instead of the reward function, we can now formulate a maximum likelihood criterion for learning a parameterized policy $\pi_\theta$. Drawing a parallel with reward modeling (see Eq.~\ref{eq:reward_model}), our policy learning objective is expressed as:

\begin{equation}\label{eq:optimum_model}
    \mathcal{L}_\text{DPO}(\pi_{\theta}; \pi_\text{ref}) = -\mathbb{E}_{(x, y_w, y_l)\sim \mathcal{D}}\left[\log \sigma \left(\beta \log \frac{\pi_{\theta}(

In this formulation, we learn an implicit reward through an alternative parameterization such that the resulting optimal policy directly corresponds to $\pi_\theta$. Additionally, since our method effectively amounts to fitting a reparameterized Bradley-Terry model, it inherits favorable theoretical guarantees, for example, consistency under appropriate conditions on the distribution of preference data \cite{bong2022generalized}. Further discussion of DPO's theoretical characteristics and comparisons to other approaches are presented in Section~\ref{sec:theory}.

\textbf{How does the DPO update function?} To gain a mechanistic insight into DPO, it is instructive to examine the gradient of its loss, $\mathcal{L}_\text{DPO}$. The gradient with respect to model parameters $\theta$ is given by:

\begin{multline*}\label{eq:gradient}
    \nabla_\theta \mathcal{L}_\text{DPO}(\pi_\theta;\pi_\text{ref}) = \\ -\beta\mathbb{E}_{(x, y_w, y_l) \sim \mathcal{D}} \bigg[\underbrace{\sigma(\hat{r}_\theta(x, y_l) - \hat{r}_\theta (x, y_w))}_\text{higher weight when reward estimate is wrong}\bigg[\underbrace{\nabla_\theta\log \pi(y_w \mid x)}_\text{increase likelihood of $y_w$} - \underbrace{\nabla_\theta\log\pi(y_l \mid x)}_\text{decrease likelihood of $y_l$}\bigg]\bigg],
\end{multline*}

where $\hat{r}_\theta(x, y) = \beta \log \frac{\pi_\theta(y \mid x)}{\pi_\text{ref}(y \mid x)}$ represents the reward implicitly modeled by the current policy $\pi_\theta$ with respect to the reference policy $\pi_\text{ref}$ (see also Section~\ref{sec:theory} for further details). Conceptually, the gradient of $\mathcal{L}_\text{DPO}$ works by encouraging higher probability for preferred continuations $y_w$ and suppressing the likelihood of less preferred $y_l$. The extent to which each sample influences the update is modulated by how much the implicit reward function $\hat{r}_\theta$ incorrectly ranks the less preferred response above the more preferred one, modulated by $\beta$—thus reflecting the severity of the reward model’s misordering and integrating the impact of the KL regularization. Empirically, this weighting scheme is crucial: omitting it, as in a simplistic version of the algorithm, leads to undesirable collapse of the language model (see Appendix Table~\ref{tab:unlikelihood_generations}).

\textbf{Overview of DPO.} The standard workflow for DPO consists of the following: First, generate two candidate completions $y_1, y_2 \sim \pi_\text{ref}(\cdot \mid x)$ for each prompt $x$, then annotate these pairs using human preference judgments to build an offline preference dataset, denoted as $\mathcal{D} = \{x^{(i)}, y_w^{(i)}, y_l)^{(i)}\}_{i=1}^N$. Second, update the language model $\pi_\theta$ by minimizing $\mathcal{L}_\text{DPO}$, using the given reference policy $\pi_\text{ref}$, dataset $\mathcal{D}$, and a selected $\beta$. 
Practically speaking, it is often desirable to utilize existing public preference datasets rather than create new completions and collect human feedback for each. Given that these datasets are often derived from $\pi^\text{SFT}$, we set $\pi_\text{ref} = \pi^\text{SFT}$ whenever this policy is accessible. In cases where $\pi^\text{SFT}$ is not obtainable, $\pi_\text{ref}$ is instead initialized by maximizing the likelihood of preferred completions ${(x, y_w)}$, specifically, by solving ${\pi_\text{ref} = \argmax_{\pi}\mathbb{E}_{x, y_w \sim \mathcal{D}}\left[\log \pi(y_w \mid x)\right]}$. This approach reduces the distribution mismatch between the ideal but unattainable reference distribution and the operational $\pi_\text{ref}$ in DPO. Further implementation specifics and chosen hyperparameters are detailed in Appendix~\ref{app:implementation}.

\section{Theoretical Analysis of DPO}
In what follows, we provide a deeper theoretical understanding of DPO, offer formal justifications, and compare the strengths of DPO to the limitations of actor-critic approaches commonly adopted in RLHF pipelines (e.g., PPO~\cite{schulman2017proximal}).

\label{sec:theory}

\subsection{Your Language Model Is Secretly a Reward Model}
DPO manages to circumvent both the need for explicit reward fitting and reinforcement learning-based policy updates by relying on a single maximum likelihood-based objective. Observe that the objective described in Eq.~\ref{eq:main_eq} corresponds to a Bradley-Terry model with reward defined as $r^*(x, y) = \beta \log\frac{\pi^*_\theta(y \mid x)}{\pi_\text{ref}(y \mid x)}$, such that optimization over the parameterized model $\pi_\theta$ is, via change of variables, equivalent to the reward model fitting objective in Eq.~\ref{eq:reward_model}. The remainder of this section develops the underlying theoretical framework for this reparameterization, demonstrates that it does not restrict the set of representable reward models, and proves that it can exactly recover the optimal policy. We begin with the formal definition of an equivalence relation among reward functions.

\begin{definition}
We define two reward functions $r(x, y)$ and $r'(x, y)$ as equivalent if and only if there exists a function $f$ such that ${r(x, y)-r'(x, y) = f(x)}$.
\end{definition}

This definition clearly constitutes an equivalence relation, leading to a partitioning of reward functions into equivalence classes. The following two lemmas can be established:

\begin{lemma}\label{lemma:same_prefrence}
Within the framework of Plackett-Luce, and specifically the Bradley-Terry model, any pair of reward functions belonging to the same equivalence class result in identical preference distributions.
\end{lemma}

\begin{lemma}\label{lemma:same_policy}
    Any two reward functions that are members of the same equivalence class will lead to the same optimal policy in the setting of a constrained reinforcement learning problem.
\end{lemma}

Proofs for these lemmas are straightforward, and are provided in Appendix \ref{app:lemma1}. The first lemma reflects the well-documented challenge of under-specification within the Plackett-Luce class of models \cite{plackett1975analysis}. This inherent ambiguity typically necessitates imposing further identifiability constraints to ensure meaningful maximum likelihood estimates as in Eq. \ref{eq:reward_model} \cite{bong2022generalized}. According to the second lemma, every reward function in a given equivalence class leads to the same optimal policy, meaning that for optimization purposes, it suffices to recover any representative from the optimal equivalence class. In Appendix~\ref{app:thm1}, we establish the following theorem:

\begin{theorem}\label{thm:main}
    Assuming only mild conditions, every reward equivalence class that aligns with the Plackett-Luce (and in particular Bradley-Terry) models can be realized through the reparameterization ${r(x, y) = \beta \log \frac{\pi(y\mid x)}{\pi_\text{ref}(y\mid x)}}$ for some model $\pi(y\mid x)$ and a fixed reference policy $\pi_\text{ref}(y \mid x)$.
\end{theorem}

\begin{sproof}
    Let $r(x, y)$ be any given reward function, which defines its associated optimal policy $\pi_r(y \mid x)$ according to Eq. \ref{eq:op_policy}. We seek to demonstrate that an equivalent reward function (from the equivalence class of $r$) can be formulated via the stated reparameterization. We introduce the projection $f$ as
\begin{equation}
    f(r; \pi_\text{ref}, \beta)(x, y) = r(x, y) - \beta\log\sum_{y}\pi_\text{ref}(y\mid x)\exp\left(\frac{1}{\beta}r(x, y)\right)
\end{equation}
Here, the mapping $f$ performs a normalization of the original reward function by subtracting the logarithm of the partition function related to $\pi_r$. Since this normalization term depends solely on $x$, $f(r; \pi_\text{ref}, \beta)(x, y)$ remains within the equivalence class of the initial $r(x, y)$. Substituting $r$ as on the right side of Eq.~\ref{eq:main_eq} (a representation that holds generally for all reward functions), we derive that $f(r; \pi_\text{ref}, \beta)(x, y) = \beta \log \frac{\pi_r(y\mid x)}{\pi_\text{ref}(y\mid x)}$. Consequently, $f$ produces a reward function within the equivalence class of $r$ with the specified structure, demonstrating that adopting this reparameterization entails no loss of generality.
\end{sproof}

An alternative interpretation of Theorem~\ref{thm:main} is that it delineates the precise reward function, within each equivalence class, that the DPO reparameterization chooses—specifically, the reward function that fulfills:

\begin{equation}\label{eq:lag_p}
     \sum_{y}\underbrace{\pi_\text{ref}(y\mid x)\exp\left(\frac{1}{\beta}r(x, y)\right)}_{=\pi(y\mid x)\text{, using Thm.~\ref{thm:main} reparam.}} = 1,
\end{equation}

which means that $\pi(y\mid x)$ constitutes a well-formed probability distribution, as it assigns non-negative values that sum to one. Notably, referencing Eq.~\ref{eq:op_policy}, Eq.~\ref{eq:lag_p} can be interpreted as the partition function associated with the optimal policy determined by the reward $r(x, y)$. The central insight of the DPO method is its ability to introduce particular constraints within the otherwise underdetermined Plackett-Luce (as well as, more specifically, the Bradley-Terry) family of preference models. This preserves the expressiveness of the space of reward models but, at the same time, guarantees that the optimal policy expressed by Eq.~\ref{eq:op_policy} remains analytically manageable for every prompt $x$.

\subsection{Instability of Actor-Critic Algorithms}  
Our framework also allows for an analysis of sources of instability in prevalent actor-critic approaches used in RLHF pipelines, such as PPO. We adopt the RLHF pipeline, with a particular emphasis on the RL fine-tuning phase as introduced in Section \ref{section:prelims}. By drawing on connections with the control as inference perspective \cite{levine2018reinforcement} for the constrained RL formulation in \ref{eq:RL}, we consider a parameterized policy model $\pi_{\theta}(y\mid x)$ and seek to minimize the divergence $\mathbb{D}_{\text{KL}}[\pi_{\theta}(y|x) \mid \mid \pi^*(y\mid x)]$, where $\pi^*$ denotes the optimal policy from Eq.~\ref{eq:optimum_model} derived from the reward function $r_{\phi}(y, x)$. After a series of algebraic manipulations, the resulting objective takes the form:

\begin{equation}\label{eq:AC}
    \max_{\pi_{\theta}}\mathbb{E}_{\pi_{\theta}(y\mid x)}\bigg[\underbrace{r_{\phi}(x, y) -\beta\log\sum_{y}\pi_\text{ref}(y\mid x)\exp\left(\frac{1}{\beta}r_{\phi}(x, y)\right)}_{f(r_{\phi}, \pi_\text{ref}, \beta)} - \underbrace{\beta\log\frac{\pi_{\theta}(y\mid x)}{\pi_\text{ref}(y\mid x)}}_{\text{KL}}\bigg]
\end{equation}

This objective aligns with that targeted in previous studies 
\citep{ziegler2020finetuning, stiennon2022learning, bai2022training, ouyang2022training}, where optimization is carried out using the DPO-equivalent reward structure corresponding to the $r_{\phi}$ reward class. In this context, the normalization factor within $f(r_{\phi}, \pi_\text{ref}, \beta)$ can be viewed as representing the soft value function associated with the reference policy $\pi_\text{ref}$. Although the presence or absence of this term does not influence the optimal solution, omitting it can cause the policy gradient estimator to exhibit high variance, potentially leading to unstable training dynamics. One solution is to model the normalization via a learned value function, though this introduces additional optimization challenges. Alternatively, a number of prior approaches employ a normalization scheme based on human completion baselines, which amounts to estimating the normalization with a single-sample Monte-Carlo method. By contrast, the DPO reparameterization constructs a reward formulation that obviates the need for such baseline-based normalization.

\section{Experiments}
Here, we present empirical studies assessing the capacity of DPO to learn policies directly from preference data. Our initial focus is on a controlled text-generation domain, examining how effectively DPO manages the trade-off between maximizing reward and minimizing KL-divergence relative to the reference policy, especially when compared to prevalent preference optimization methods like PPO. Subsequently, we test DPO on larger-scale models and more demanding RLHF benchmarks, such as tasks in summarization and dialogue. Across these evaluations, DPO generally matches or exceeds the performance of strong baselines like RLHF with PPO, frequently excelling even without significant hyperparameter tuning. Additionally, it often surpasses the strategy of selecting the top-performing trajectory from $N$ samples under a learned reward. Prior to detailing the outcomes, we outline the experimental methodology; further specifics are available in Appendix~\ref{app:exp_details}.

\textbf{Tasks.} We investigate three distinct open-ended text generation scenarios. Across all settings, each method optimizes a policy from a preference dataset $\mathcal{D}=\bigl\{x^{(i)}, y_w^{(i)}, y_l^{(i)}\bigr\}_{i=1}^N$. For the \textbf{controlled sentiment generation} task, $x$ is a prompt derived from a movie review in the IMDb dataset \cite{maas-EtAl:2011:ACL-HLT2011}, and the objective is to generate a continuation $y$ exhibiting positive sentiment. To ensure systematic evaluation, we automatically construct preference pairs for each generated output using a pre-trained sentiment classifier, selecting pairs such that $p(\text{positive}\mid x,y_w)>p(\text{positive}\mid x,y_l)$. In this setup, SFT involves fine-tuning GPT-2-large on the training partition of IMDb movie reviews until performance plateaus (additional procedural details in App~\ref{app:sentiment_details}). For the \textbf{summarization} task, $x$ corresponds to a Reddit forum post, while the policy is trained to produce a summary $y$ that concisely distills the main ideas of the post. Building on prior approaches, we utilize the Reddit TL;DR dataset \citep{volske-etal-2017-tl} supplemented with human-annotated preferences as collected by \citeauthor{stiennon2022learning}. The SFT baseline is adapted from a model trained on human summaries of forum posts\footnote{\url{https://huggingface.co/CarperAI/openai_summarize_tldr_sft}} and is used within the TRLX framework \citep{leandro_von_werra_2023_7790115} to facilitate RLHF training. The preference annotations used were produced by \citeauthor{stiennon2022learning} from samples generated by a related, yet independently-trained, SFT model. In the \textbf{single-turn dialogue} setting, $x$ is a user-submitted question or prompt that might range from a scientific inquiry to a request for personal guidance. The policy must formulate a helpful and engaging response $y$. For this task, we leverage the Anthropic Helpful and Harmless dialogue corpus \citep{bai2022training}, which contains 170,000 conversations between a human user and an automated assistant. Each dialogue concludes with two assistant-generated responses and an associated label identifying the preferred answer, as judged by a human annotator. Since a suitable pre-trained SFT model is not available in this context, we create the SFT baseline by fine-tuning a generic language model solely on the preferred completions within this dataset.

\textbf{Evaluation.} We adopt two distinct evaluation methodologies in our experiments. To specifically assess how well each algorithm solves the constrained reward maximization problem, we use the controlled sentiment generation scenario, measuring each method by the set of reward and KL-divergence pairs it achieves relative to the reference policy. This is feasible in this setting because we possess the true reward function in the form of a sentiment classifier. Conversely, outside this controlled environment where the actual reward function remains unknown, we resort to measuring an algorithm's \textit{win rate} when compared to a baseline policy. For this, we utilize GPT-4 to serve as a stand-in for human judges, evaluating summary quality for the summarization task and helpfulness of responses in the single-turn dialogue task. Summarization tasks use ground-truth summaries from the test set as the baseline comparison; in dialogue, the test set’s preferred responses are treated as the baseline. Recent literature indicates that large language models can offer superior automatic evaluations compared to classic metrics \citep{Chen2023ExploringTU}. However, in Section~\ref{sec:human-judgments}, we supplement this with a human subject study to validate our deployment of GPT-4 for evaluation. Our results show that GPT-4's evaluations exhibit a strong correlation with human judgments, with rates of agreement comparable to, or exceeding, that observed among human annotators.

\textbf{Methods.} Beyond DPO, we benchmark a range of prevailing strategies for aligning language model outputs with human preference signals. For baseline prompting approaches, we employ zero-shot prompting with \textbf{GPT-J} \citep{gpt-j} in summarization and two-shot prompting using \textbf{Pythia-2.8B} \citep{biderman2023pythia} in the dialogue setting. Additionally, we assess the performance of the standard \textbf{SFT} model and introduce \textbf{Preferred-FT}—a model trained via supervised fine-tuning on the selected completion $y_w$. This is done by drawing $y_w$ from the SFT model in the sentiment and summarization settings, and from a general-purpose LM in dialogue. Another approach we consider is the pseudo-supervised \textbf{Unlikelihood} method~\citep{welleck2019neural}, which encourages the model to increase the probability of $y_w$ while actively reducing the likelihood of $y_l$, modulated by a hyperparameter $\alpha\in[0,1]$ for the unlikelihood loss. Further, we include \textbf{PPO} \citep{schulman2017proximal}, utilizing a reward model trained on preference data, as well as \textbf{PPO-GT}—an oracle setup where the policy is trained with access to the ground-truth reward in the sentiment domain. In sentiment experiments, PPO-GT is implemented in two forms: one uses an out-of-the-box version \cite{leandro_von_werra_2023_7790115}, and the other involves a version with reward normalization and fine-tuned hyperparameters for optimized results (these same modifications are applied to standard PPO with learned rewards). As a final comparator, we introduce the \textbf{Best of $N$} approach: $N$ samples are generated from the SFT (or Preferred-FT in dialogue) and the sample with the top score under a reward model trained on preference data is chosen. Although this strategy can deliver strong results, it separates reward model quality from PPO optimization and is limited in practicality due to the need to sample $N$ completions for every query during evaluation, making it computationally intensive even for modest values of $N$.

\subsection{How well can DPO optimize the RLHF objective?}

In conventional RLHF methods, the objective function maximizes the reward while imposing a KL constraint to limit divergence from the reference policy. Thus, algorithm comparisons must jointly consider both the reward obtained and the corresponding KL distance; attaining a marginally higher reward at the expense of a substantially larger KL is not advantageous. Figure~\ref{fig:frontier-tldr-main} presents the trade-off curve between reward and KL for different algorithms applied to the sentiment task. For each method, we conduct several training sessions, varying the policy's conservativeness parameter (target KL $\in\{3,6,9,12\}$ for PPO, $\beta \in \{0.05,0.1,1,5\}$, $\alpha\in\{0.05,0.1,0.5,1\}$ for unlikelihood, and random seeds for preferred-FT), for a total of 22 trials. Every 100 training updates up to convergence, we assess each learned policy using a suite of test prompts. The evaluations report the mean reward based on the true reward model and the average sequence-level KL\footnote{Specifically, the aggregate KL divergence across all time steps.} with respect to the reference policy, $\text{KL}\left(\pi\mid \mid \pi_\text{ref}\right)$. Our findings indicate that DPO establishes a reward-KL frontier that is substantially more efficient than competing methods, achieving superior rewards without incurring excessive KL penalties. Several aspects of this result are particularly noteworthy. While both DPO and PPO target the same learning objective, DPO consistently yields a more favorable reward/KL balance, outperforming PPO along this frontier. Additionally, DPO surpasses PPO even in conditions where PPO benefits from access to ground truth rewards (PPO-GT).

\subsection{Can DPO scale to real preference datasets?}
\label{sec:dpo-real-datasets}
We proceed by assessing the fine-tuning capabilities of DPO within the domains of summarization and single-turn dialogue. In summarization, widely used automatic metrics like ROUGE have shown weak alignment with human preferences~\citep{stiennon2022learning}, and earlier research demonstrates that fine-tuning language models using PPO guided by human feedback can yield more effective summaries. Our experimental setup involves generating completions on the test partition of the TL;DR summarization dataset and measuring the mean win rate versus the reference completions from the test set. For all approaches, outputs are generated across a temperature spectrum from 0.0 to 1.0, with the corresponding win rates illustrated in Figure~\ref{fig:frontier-tldr-main} (right). DPO, PPO, and Preferred-FT are all initialized from the identical GPT-J SFT model\footnote{\url{https://huggingface.co/CarperAI/openai_summarize_tldr_sft}}. Results show that at temperature 0.0, DPO reaches an approximate win rate of 61\%, surpassing PPO, which attains around 57\% at its optimal temperature setting. Furthermore, DPO achieves a superior peak win rate relative to the best among the $N$ baseline models. It is important to mention that the $\beta$ hyperparameter in DPO was not thoroughly optimized, indicating these figures may not represent the method’s upper bound. Additionally, DPO displays significantly greater resilience to variations in sampling temperature compared to PPO, whose effectiveness deteriorates at higher temperatures to levels similar to the base GPT-J model. Improvements from Preferred-FT over the SFT baseline are marginal. Direct human evaluation comparisons between DPO and PPO are discussed in Section~\ref{sec:human-judgments}, where completions generated by DPO at temperature 0.25 were favored 58\% of the time relative to PPO outputs at temperature 0.

For single-turn dialogue tasks, we assess the various approaches using the subset of the Anthropic HH dataset's test split \citep{bai2022training} that contains a single exchange between human and assistant. To quantify performance, we use GPT-4 to determine win rates for each approach, referencing the completions marked as preferred in the test data. Due to the absence of a standard SFT model for this benchmark, our pipeline begins with the pre-trained Pythia-2.8B model. We fine-tune a reference model with Preferred-FT, aligning its output distribution with the chosen completions, and subsequently train a model via DPO. For further comparison, we evaluate against the best response out of 128 Preferred-FT completions—establishing 128 as the saturation point for the Best of $N$ method for this setting (see Appendix Figure~\ref{fig:best-of-n})—and employ a 2-shot prompt using the Pythia-2.8B base model. We observe that DPO consistently matches or exceeds the best performances from each method's optimal temperature. Additionally, an RLHF model trained via PPO on the Anthropic HH dataset\footnote{\url{https://huggingface.co/reciprocate/ppo_hh_pythia-6B}} (sourced from a widely recognized implementation\footnote{\url{https://github.com/CarperAI/trlx/tree/main/examples/hh}}) is assessed, but no combination of prompting or temperature produces superior results to the original Pythia-2.8B model. Drawing on our TL;DR findings and the shared reward objective between methods, we treat Best of 128 as an approximate surrogate for PPO-level outcomes. Overall, DPO emerges as the only method that not only remains computationally efficient but also improves over the preferred completions within the Anthropic HH dataset, matching or exceeding the performance of the resource-intensive Best of 128 strategy. Notably, Figure~\ref{fig:dialogue-main} indicates that DPO attains its peak performance after relatively few training steps.

\subsection{Generalization to a new input distribution}

To deepen our analysis of how PPO and DPO respond to changes in input distribution, we tested the trained policies from the Reddit TL;DR summarization task on a novel data distribution: the test set of the CNN/DailyMail news articles dataset \citep{nallapati-etal-2016-abstractive}. The policies were sampled using the optimal temperatures identified in the TL;DR setting (0 and 0.25, respectively). Table~\ref{tab:ood} summarizes the outcomes. Using the same GPT-4 (C) prompt as in the Reddit evaluation—except that ``forum post'' was replaced with ``news article''—we assessed GPT-4's win rate against the reference summaries from the CNN/DailyMail dataset. Notably, DPO consistently achieves a considerable advantage over PPO under this distribution shift. These findings indicate that DPO policies are capable of generalizing at least as effectively as those obtained with PPO, despite DPO's lack of access to the supplementary unlabeled Reddit TL;DR prompts that PPO leverages.

\subsection{Validating GPT-4 judgments with human judgments}
\label{sec:human-judgments}
To assess the dependability of GPT-4 as an evaluation tool, we ran a human subject study using summaries from our TL;DR experiments and two different GPT-4 prompts. The \textbf{GPT-4 (S)} (simple) prompt asks directly which summary better conveys the main information from the post. In contrast, the \textbf{GPT-4 (C)} (concise) prompt instructs GPT-4 to also consider conciseness, motivated by our observation that the simple prompt leads GPT-4 to favor lengthier and more repetitive outputs relative to human raters. Complete details of both prompts are provided in Appendix~\ref{app:prompts}. Our human evaluation encompassed three main system comparisons—one with the best performing model (DPO, temp. 0.25), one with the least effective (PPO, temp. 1.0), and one of intermediate performance (SFT, temp. 0.25)—each evaluated against PPO using greedy sampling, in order to reflect a range of output qualities. Human judgments were collected for all DPO and PPO-1 matchups, though due to the small pool of human raters, we gathered multiple human ratings only for these pairs. Results reveal that for both prompts, the level of agreement between GPT-4 and human annotators matches the inter-human agreement rate, supporting the use of GPT-4 as a stand-in for human assessments. Furthermore, the \textbf{GPT-4 (C)} prompt tends to yield win rates that more closely align with human preferences; accordingly, we use this prompt in reporting primary results in Section~\ref{sec:dpo-real-datasets}. Further methodological information, including details of the rating interface and the identities of participating annotators, is available in Appendix~\ref{app:human-study}.

\section{Discussion}
Preference-based learning constitutes a robust and scalable methodology for developing language models that are both competent and well-aligned with human values. In this work, we have presented DPO, a straightforward approach that enables training language models using preference data without the need for reinforcement learning. Instead of adapting the preference learning task to fit conventional RL frameworks and thereby utilize standard RL tools, DPO leverages an explicit correspondence between language model policies and reward functions. This allows for the direct training of language models to align with human preferences by employing a standard cross-entropy loss, eliminating the requirement for reinforcement learning and avoiding any compromise in generality. Notably, with almost no hyperparameter adjustment, DPO achieves performance comparable to, or exceeding, state-of-the-art RLHF methods, including those relying on PPO. Consequently, DPO significantly lowers the threshold for training language models using human feedback.

\textbf{Limitations \& Future Work.} The findings presented here open up several avenues for subsequent research. An open question is the extent to which the DPO-trained policies are able to generalize beyond the training distribution, especially when compared to models trained with explicit reward signals. Preliminary observations indicate DPO policies may exhibit generalization capabilities similar to those obtained with PPO, though a thorough investigation remains to be conducted. For example, an interesting line of inquiry is whether self-labeling through the DPO framework can effectively exploit prompts lacking labels. Another topic concerns the characteristics and implications of reward over-optimization within the direct preference optimization context—whether the marginal reduction in performance shown in Figure~\ref{fig:dialogue-main}-right exemplifies this phenomenon. Furthermore, while current experiments consider models up to 6B parameters, scaling DPO to models with far greater capacity presents a promising opportunity for future exploration. In terms of model assessment, our analysis reveals that win rates produced by GPT-4 are sensitive to prompt phrasing, suggesting further study is warranted into optimizing the quality of automated system evaluations. Ultimately, the potential applications of DPO extend well beyond language models trained on human preferences, offering promise for training generative systems in various domains.